% Capítulo 2 - Fundamentação Teórica

\chapter{Fundamentação Teórica}
\label{cap:fundamentacao}

\section{Conceito e Função de um NAS Doméstico}

NAS é a sigla para \textit{Network Attached Storage}, isto é, armazenamento conectado à rede. Em termos práticos, um NAS é um servidor dedicado a centralizar arquivos e disponibilizá-los por protocolos de rede. Em vez de conectar um disco diretamente a um único computador, o NAS permite que diferentes dispositivos da residência acessem os mesmos dados pela rede local.

Em ambientes domésticos, o NAS pode funcionar como repositório central de documentos, destino de backup, armazenamento de fotos e vídeos, compartilhamento local e base para serviços complementares. Para fins acadêmicos, o NAS é um objeto adequado porque permite medições de \textit{throughput}, taxa de leitura e escrita, custo por capacidade, consumo energético e comportamento sob carga.

\section{Armazenamento Local, Nuvem Pública e Nuvem Privada}

Armazenamento local e nuvem pública não devem ser tratados como categorias excludentes. O armazenamento local privilegia autonomia, proximidade física dos dados e menor dependência de terceiros. A nuvem pública privilegia conveniência, escalabilidade e integração entre dispositivos. A nuvem privada ocupa uma posição intermediária: os dados permanecem sob controle do usuário, mas são acessados por uma camada de serviços web, sincronização e colaboração.

\section{Single Board Computers como Base de Infraestrutura Doméstica}

Single Board Computers são computadores completos concentrados em uma única placa, geralmente com CPU, memória, interfaces USB, rede e armazenamento de \textit{boot}. Plataformas ARM de baixo custo são utilizadas em ensino, automação, IoT, prototipagem e servidores leves. Suas vantagens incluem baixo consumo energético, formato compacto e compatibilidade com Linux.

A Radxa ROCK 4B+ será utilizada como plataforma de bancada para avaliar a viabilidade de um NAS doméstico com três SSDs externos. A escolha por SSDs conectados por adaptador externo USB-SATA reduz custo e simplifica a montagem, mas também cria pontos de atenção relacionados a energia, barramento USB, aquecimento e estabilidade.

\section{OpenMediaVault}

O OpenMediaVault é apresentado em sua documentação oficial como uma solução completa de NAS baseada em Debian Linux, com interface web de administração e suporte a protocolos como FTP, SMB/CIFS e NFS. No projeto, ele atua como camada de gerenciamento do NAS, centralizando montagem de discos, criação de pastas compartilhadas, usuários, permissões e serviços de arquivo.

\section{Nextcloud}

De acordo com a documentação oficial do Nextcloud~\cite{nextcloud2026}, é uma plataforma \textit{open source} de sincronização, compartilhamento de arquivos e colaboração \textit{self-hosted}. No contexto do projeto, ele não substitui o NAS; ele opera sobre o armazenamento administrado pelo sistema, oferecendo experiência de nuvem privada com interface web, clientes desktop e móveis, \textit{upload}, \textit{download}, compartilhamento e sincronização.

\section{Tailscale e Acesso Remoto Seguro}

A exposição direta de serviços domésticos à internet aumenta a superfície de ataque. O Tailscale oferece uma alternativa prática ao criar uma \textit{mesh VPN} baseada em WireGuard~\cite{tailscale2026}, com autenticação e controle de dispositivos autorizados. Neste projeto, ele reduz a necessidade de abertura direta de portas no roteador e permite testes remotos de acesso ao Nextcloud e, quando adequado, aos serviços administrativos.

\section{Pontos de Medição e Limitações Esperadas}

A adoção de uma SBC com SSDs conectados por adaptador externo USB-SATA impõe limitações que devem ser observadas. Entre elas estão possíveis gargalos no barramento USB, variação de qualidade dos adaptadores, demanda elétrica dos SSDs, aquecimento, estabilidade sob carga e desempenho inferior ao de NAS comerciais com controladores SATA dedicados. Essas limitações não invalidam o projeto; elas definem justamente o que precisa ser medido.

\begin{figure}[htb]
    \legenda[fig:gargalos]{Pontos de medição e possíveis gargalos}
    \fig{width=0.9\textwidth}{assets/img/gargalos}
    \fonteautoria
\end{figure}